% ===================================================================
%  TABEL Nr. 14 — Tänav. --- Kaupmehed.
%  Traduction 2026 d'apres texto/io, controlee sur texto/fr et
%  texto/fi. Consigne : voir texto/et/10-tabel-01.tex.
%
%  LE TABLEAU 14 LITUANIEN AVAIT TROUVE UNE LIGNE SUR LA MEME RUE ;
%  L'ESTONIEN LA RETROUVE, ET COMME LES DEUX LANGUES N'ONT AUCUNE
%  PARENTE, CE N'EST PLUS UNE OBSERVATION MAIS UNE REGLE.
%
%  Le lituanien nommait soixante metiers avec quatre suffixes, et les
%  seules exceptions etaient les diplomes. L'estonien fait le meme
%  partage avec les siens :
%
%    « -ja » sur un verbe : « müüja » le vendeur, de « müüma » ;
%      « ostja » l'acheteur, « kandja » le porteur, « pakkija »
%      l'emballeur, « pühkija » le balayeur, « köitja » le relieur,
%      « kirjastaja » l'editeur, « trükkal » l'imprimeur ;
%    « -nik » sur un nom : « aednik » le jardinier, « kunstnik »
%      l'artiste, « ametnik » l'employe ;
%    « -sepp » du tableau 5 : « plekksepp » le ferblantier,
%      « kellassepp » l'horloger, « kübarasepp » le chapelier.
%
%  ET LES EXCEPTIONS SONT LES MEMES QU'A VILNIUS : « advokaat »,
%  « notar », « insener », « arhitekt », « hambaarst », « optik »,
%  « fotograaf », « litograaf ».
%
%  LA LIGNE N'EST PAS ENTRE L'ANCIEN ET LE NOUVEAU. L'imprimeur est
%  « trükkal », mot estonien, et son metier a quatre siecles ; le
%  photographe est « fotograaf », mot international, et le sien en a
%  cinquante. LA LIGNE EST ENTRE CE QU'ON APPREND DANS UN ATELIER ET
%  CE QU'ON APPREND POUR UN EXAMEN. Un apprenti relieur regarde son
%  maitre ; un avocat passe devant un jury, et le jury, la faculte,
%  le titre et le mot lui viennent tous ensemble du meme dehors.
%
%  Deux langues sans rapport, deux systemes de suffixes sans rapport,
%  la meme frontiere au meme endroit d'une meme rue : c'est la
%  troisieme fois que la serie verifie une regle sur deux familles a
%  la fois, apres la caisse d'epargne du tableau 11 romanche et le
%  menuisier du tableau 5 estonien.
% ===================================================================

%%K t14-apar-1 apar
\VUsaut{4.0mm}
\VUtitre{15.0pt}{60}{TABEL Nr. 14}
\VUsaut{3.0mm}

%%K t14-tit-1 sub
\VUpk{11.4pt}{40}{Tänav. --- Kaupmehed.}
\VUsaut{3.0mm}

%%K t14-01-1 p
1. --- Mu sõber Jaan, kes elab maal, tuli veetma kolm päeva minu perekonnas.
\VUgras{Seni} ei olnud tal olnud juhust näha suurt linna; seepärast veedame me kõik oma
päevad jalutades, ja ma seletan talle seda, mida ta ei tunne.

%%K t14-02-1 p
2. --- Kõigepealt läksime me külastama \VUgras{observatooriumi} \textsuperscript{(1)},
mis teda väga huvitas. Tagasi tulles \VUgras{tänava} \textsuperscript{(3)} kaudu, kus on
\VUgras{türgi saun} \textsuperscript{(2)}, kohtasime me \VUgras{matuserongi}
\textsuperscript{(4)}. \VUgras{Matusesaatja} ees kõndisid \VUgras{vaimulikud}, kes
eelnesid \VUgras{surnuvankrile}, millesse oli pandud \VUgras{kirst}. Palju sõpru saatis
\VUgras{leinavat} perekonda.

%%K t14-02-2 p
Pärast rongkäigu möödumist ületasime me tänava suure \VUgras{õllesaali}
\textsuperscript{(5)} ees, kus paljud \VUgras{külastajad} jõid \VUgras{õlut} terrassil,
varjatuna klaasist \VUgras{markiisiga} \textsuperscript{(6)}.

%%K t14-03-1 p
3. --- Jaan oli väga hämmastunud \VUgras{vapist}, mis oli pandud
\VUgras{likööripoe} \textsuperscript{(7)} ja \VUgras{muusikakaupmehe}
\textsuperscript{(10)} kaupluse vahele. Ta küsis minult selle tähendust; ma vastasin, et
see näitab Inglise \VUgras{konsulaati} \textsuperscript{(8)}, ja juhtisin ta tähelepanu
\VUgras{konsulile} \textsuperscript{(9)}, kes on rõdul.

%%K t14-tit-2 sub
\VUpk{10.2pt}{40}{Poodides vaatamine.}
\VUsaut{3.0mm}

%%K t14-04-1 p
4. --- Muidugi peatusime me \VUgras{muusikariistade}, \VUgras{harmooniumide},
\VUgras{orelite} ja \VUgras{klaverite} väljapaneku ees; me vaatasime klaveri
\VUgras{nootide} ja \VUgras{muusikapalade} illustreeritud kaasi ja imetlesime kullatud
puust \VUgras{lüürat}, mida kasutatakse sildina.

%%K t14-05-1 p
5. --- Tolmupilved, mille tegid \VUgras{tänavapühkijad} \textsuperscript{(11)} ja
\VUgras{pühkimisvanker} \textsuperscript{(12)}, sundisid meid kiiresti mööduma
\VUgras{mööblikaupmehe} \textsuperscript{(13)} ilusast \VUgras{mööblist} ja
\VUgras{plekksepa} \textsuperscript{(14)} \VUgras{ahjude} ja \VUgras{lampide} väljapanekust.

%%K t14-06-1 p
6. --- Pealegi olime me siin kohal sunnitud kõnniteelt lahkuma, sest see oli täis
pakke, mida võeti \VUgras{veovankrilt} \textsuperscript{(16)}, et panna neid
\VUgras{kauplusse} \textsuperscript{(15)}, mis oli äsja üüritud.

%%K t14-06-2 p
See takistas meil nägemast \VUgras{vankrimeistri} \textsuperscript{(17)} ilusaid vankreid,
kuid ei takistanud meid peatumast, et vaadata \VUgras{pakkijat} \textsuperscript{(19)},
kes tegi kasti oma naabrile, \VUgras{jalgrataste ja autode müüjale}
\textsuperscript{(18)}. Siis, kuna oli juba veidi hilja, läksime me koju tagasi.

%%K t14-tit-3 sub
\VUpk{10.2pt}{40}{Jalutuskäik linnas.}
\VUsaut{3.0mm}

%%K t14-07-1 p
7. --- Järgmisel päeval tegi mu ema Jaanile ettepaneku, et teda pildistataks, siis
tegi ta mulle ülesandeks teda saata \VUgras{fotograafi} \textsuperscript{(20)} juurde.
Pärast lõunasööki läksime me kunstniku ateljeesse, kus me olime sunnitud üsna kaua
ootama. Et end tegevuses hoida, vaatasime me \VUgras{portreesid} \textsuperscript{(22)},
mis seinal ripuvad.

%%K t14-07-2 p
Kui tuli meie kord, ei kestnud töö kaua, ja niipea kui mu sõber lõpetas poseerimise
\VUgras{fotoaparaadi} \textsuperscript{(21)} ees, läksime me alla.

%%K t14-08-1 p
8. --- Teisel korrusel nägime me läbi \VUgras{juuksuri} \textsuperscript{(23)} ukse, mis
oli avatud, tema abilist kliendi juukseid lõikamas.

%%K t14-08-2 p
Tänavale jõudnud, möödusime me \VUgras{tubakapoest} \textsuperscript{(24)}. Jaan oli nii
lõbustatud punasest \VUgras{laternast} \textsuperscript{(25)} ja \VUgras{suurest piibust},
mida kasutatakse \VUgras{sildina} \textsuperscript{(26)}, et ta põrkas kokku kahe vana
härraga, kes vestlesid \VUgras{nuusktubakat nuusates} \textsuperscript{(27)}.

%%K t14-tit-4 sub
\VUpk{10.2pt}{40}{Kondiitriäri.}
\VUsaut{3.0mm}

%%K t14-09-1 p
9. --- Kõigi maade \VUgras{pangatähed} ja \VUgras{mündid}, mis olid välja pandud
\VUgras{rahavahetaja} \textsuperscript{(29)} \VUgras{vitriini}, köitsid mõneks hetkeks
meie tähelepanu, kuid kuna oli oode aeg, astusime me \VUgras{kondiitri}
\textsuperscript{(31)} juurde kauem peatumata.

%%K t14-10-1 p
10. --- Nähes kõiki \VUgras{maiustusi}, mida kauplus sisaldas, \VUgras{kooke,
piparkooke, küpsiseid} ja igasuguseid \VUgras{kompvekke}, ei suutnud me kergesti valida.
Lõpuks valis Jaan \VUgras{koorekoogid}, ja mina võtsin ainult väikese
\VUgras{kirsitordi}, sest maiustused \VUgras{teevad mu hambad haigeks} ja ma ei taha
sugugi liiga sageli käia \VUgras{hambaarsti} \textsuperscript{(30)} juures.

%%K t14-tit-5 sub
\VUpk{10.2pt}{40}{Masinal kirjutamine.}
\VUsaut{3.0mm}

%%K t14-11-1 p
11. --- Et näidata oma sõbrale \VUgras{kirjutusmasinat}, millest ta oli kuulnud
rääkivat, kuid mida ta ei tundnud, astusime me \VUgras{kontorisse} \textsuperscript{(32)},
mis kuulub minu \VUgras{nõole} \textsuperscript{(34)}. Me leidsime ta oma \VUgras{sekretärile}
\textsuperscript{(35)} kirju dikteerimas, kes on \VUgras{stenografist}. Vaevalt oli kiri
lõpetatud, kui \VUgras{masinakirjutaja} \textsuperscript{(33)} selle oma masinal ümber
kirjutas. Soovimata katkestada oma sugulase tegevust, jäime me tema juurde ainult
mõneks hetkeks.

%%K t14-12-1 p
12. --- Mu nõo kontori all elab suur \VUgras{kellassepp-juveliir}
\textsuperscript{(37)}, kes on pealegi kullassepp. Tema suurepärane kauplus sisaldab
palju \VUgras{kellasid, pendelkellasid, taskukellasid, ketikesi} ja hinnalisi
\VUgras{ehteid}, näiteks \VUgras{pärlikeesid}, kuldseid \VUgras{käevõrusid},
\VUgras{kõrvarõngaid} ja \VUgras{sõrmuseid}, mis on ehitud \VUgras{kalliskivide} ja
\VUgras{teemantidega}. Üks vitriinidest on eriliselt reserveeritud \VUgras{kuldesemetele}
ja sisaldab tähtsat hõbedaste \VUgras{söögiriistade} kogu.

%%K t14-13-1 p
13. --- Tänavanurgal keerates märkasin ma, et oli vahetatud \VUgras{silt}
\textsuperscript{(36)}, mis on pandud maja \VUgras{numbri} juurde. Ma olin oma märkust
valjusti ütlemas, kui kuuldus \VUgras{sõjaväeorkestri} rõõmsaid helisid. Me jooksime
rügemendile vastu, mõtlemata, et \VUgras{see} pidi mööduma \VUgras{kübarasepa}
\textsuperscript{(39)} ja \VUgras{kindakaupmehe} \textsuperscript{(40)} kaupluste eest ja
et me oleksime olnud sama hästi paigutatud, et selle paraadi näha, jäädes
\VUgras{optiku} \textsuperscript{(38)} vitriini ette, mis on täis \VUgras{näpitsprille,
lornette} ja \VUgras{pikksilmi}.

%%K t14-tit-6 sub
\VUpk{10.2pt}{40}{Marsiparaad.}
\VUsaut{3.0mm}

%%K t14-14-1 p
14. --- \VUgras{Rügemendi} \textsuperscript{(41)} ees kõndis \VUgras{trummimajor}
\textsuperscript{(42)}, kellele järgnesid \VUgras{trummarid} \textsuperscript{(43)},
\VUgras{signaalpasunad} \textsuperscript{(44)} ja kogu \VUgras{orkester}.
\VUgras{Kindral} \textsuperscript{(45)}, kes oli platsil oma adjutandiga, oli peatunud
\VUgras{veejoa} \textsuperscript{(49)} juures, mis tuleb \VUgras{purskkaevust}
\textsuperscript{(48)}, et \VUgras{paraadi} vaadata.

%%K t14-14-2 p
\VUgras{Staabiohvitserid} \textsuperscript{(46)}, \VUgras{kolonel} ja
\VUgras{kolonelleitnant} tervitasid teda mõõgaga. \VUgras{Majorid} olid oma
\VUgras{pataljonide} \textsuperscript{(47)} ees ja iga \VUgras{kapten} juhtis oma
\VUgras{kompaniid}, sellal kui \VUgras{leitnandid}, \VUgras{nooremleitnandid} ja
\VUgras{allohvitserid} olid oma tavalisel kohal oma meeste kõrval.

%%K t14-15-1 p
15. --- Palju möödujaid oli kogunenud \VUgras{tänavaristile} \textsuperscript{(50)};
kaupmehed, \VUgras{vihmavarjumeister} \textsuperscript{(51)}, \VUgras{värvija}
\textsuperscript{(53)}, \VUgras{relvasepp} \textsuperscript{(54)} tulid välja
\VUgras{sõdureid} vaatama. Kõik sõidukid, \VUgras{voorimehed} \textsuperscript{(52)},
\VUgras{mastvankrid} \textsuperscript{(57)} ja ka \VUgras{kastmisvaat}
\textsuperscript{(56)}, seadsid end kõnnitee äärde.

%%K t14-tit-7 sub
\VUpk{10.2pt}{40}{Stseenid tänaval.}
\VUsaut{3.0mm}

%%K t14-15-2 p
\VUgras{Kaubareisija} \textsuperscript{(55)}, kes kandis käes oma
\VUgras{näidiskaste}, ületas kiiresti tänava, ja inimesed, kes istusid
\VUgras{pealmisel korrusel} \textsuperscript{(59)}, mis on \VUgras{omnibussil}
\textsuperscript{(58)}, pöördusid ümber, et vaatepilti nautida. Vaene \VUgras{rändlaulja}
\textsuperscript{(60)} oma \VUgras{leierkastiga} \textsuperscript{(61)} pidi katkestama oma
\VUgras{laulu}, sest trummide lärm kattis tema hääle.

%%K t14-16-1 p
16. --- Kui rügement jõudis \VUgras{riidekaupluse} \textsuperscript{(64)} ette, jätsid
\VUgras{õmblejannad} \textsuperscript{(63)} ja \VUgras{rätsepad} \textsuperscript{(62)} oma
\VUgras{õmblusmasinad} ja oma töö, et aknast vaadata. \VUgras{Kaupmees}
\textsuperscript{(65)}, kes oli äsja pannud uued \VUgras{ülikonnad}
\textsuperscript{(66)} oma väljapanekusse, pidi laskma sisse viia \VUgras{kalevid}
\textsuperscript{(67)} ja \VUgras{linased}, mis olid seatud kõnniteele
\VUgras{kanalisatsiooniluugi} \textsuperscript{(68)} juurde, kartes, et rahvahulk need
ümber ajab; isegi vana \VUgras{rändkaupmees} \textsuperscript{(69)}, kellelt üks portjee
sukkade eest tingis, oli sunnitud koridori astuma, et oma müük lõpetada.

%%K t14-16-2 p
Me saatsime rügementi kasarmuni \VUgras{ja,} oma päevaga väga rahul, tulime me tagasi
õhtusöögi ajaks.

%%K t14-tit-8 sub
\VUpk{10.2pt}{40}{Mitmesugused kaubad ja ametid.}
\VUsaut{3.0mm}

%%K t14-17-1 p
17. --- Järgmisel päeval tegi mu isa meile ettepaneku külastada sealse ajalehe
trükikoda. Me võtsime selle vaimustusega vastu.

%%K t14-17-2 p
\VUgras{Trükkal} on ka \VUgras{raamatukaupmees} \textsuperscript{(70)}
\VUgras{(= raamatute müüja)}, \VUgras{kirjastaja}, \VUgras{paberikaupmees} ja
\VUgras{köitja}. Ta on seadnud teisele korrusele ajalehe \VUgras{toimetuse}
\textsuperscript{(71)} ja ka \VUgras{graveerimistoa}, kus töötavad \VUgras{litograafid}
\textsuperscript{(72)} ja \VUgras{vaskgravöörid} \textsuperscript{(73)}, lõpuks
\VUgras{laduruumi}, kus on \VUgras{trükitöölised} \textsuperscript{(74)}. Päris
\VUgras{trükikoda} \textsuperscript{(75)}, \VUgras{trükipressid} ja mitmesugused masinad
on esimesel korrusel.

%%K t14-tit-9 sub
\VUpk{10.2pt}{40}{Jaan esitab väga palju küsimusi.}
\VUsaut{3.0mm}

%%K t14-18-1 p
18. --- Trükikojast väljudes peatus Jaan väga hämmastununa väikese klaaskarbi ees, mis
oli pandud sambale, vastu \VUgras{pudukauplust} \textsuperscript{(77)}, ja küsis, milleks
seda kasutatakse. Mu isa seletas talle, et see on \VUgras{tulekahjuteataja}
\textsuperscript{(76)}. Pealegi oli linnas kõik mu sõbra jaoks hämmastav, alates lihtsast
\VUgras{kivipurskkaevust} \textsuperscript{(78)} ja kergest \VUgras{veokolmerattalisest}
\textsuperscript{(80)}, mida juhtis livreega \VUgras{pikkolo} \textsuperscript{(79)}, kuni
\VUgras{postivankrini} \textsuperscript{(81)}.

%%K t14-19-1 p
19. --- Sellal kui mu isa meid mõneks hetkeks jättis, et rääkida \VUgras{maksukogujaga}
\textsuperscript{(83)}, kes oli peatunud vestlema \VUgras{advokaadi}
\textsuperscript{(82)} ja \VUgras{notariga} \textsuperscript{(84)}, lõbustasime me end,
jälgides silmadega liiklust tänaval. Meist möödusid kõige mitmekesisemad inimesed:
\VUgras{kondiitri poiss} \textsuperscript{(85)}, kes kandis korvis suurepärast
\VUgras{pasteeti}, tuli \VUgras{rõivakaupmehe} \textsuperscript{(86)} järel;
\VUgras{laternasüütaja} \textsuperscript{(87)} vestles \VUgras{gaasitöölisega}
\textsuperscript{(88)}; \VUgras{nunn} \textsuperscript{(89)}, kes hoidis käest orvukest,
läks tagasi oma kloostrisse.

%%K t14-tit-10 sub
\VUpk{10.2pt}{40}{Meie koduteel.}
\VUsaut{3.0mm}

%%K t14-20-1 p
20. --- Hetkel, mil mu isa meiega uuesti ühines, naeris Jaan valjusti, nähes kahe
tahmase \VUgras{korstnapühkija} \textsuperscript{(91)} musta nägu ja kummalist riietust.

%%K t14-21-1 p
21. --- Ma tahtsin osta \VUgras{lillemüüjalt} \textsuperscript{(92)} taime oma emale,
kuid mu isa juhtis mu tähelepanu sellele, et see oleks kaasakandmiseks liiga raske ja
et parem oleks osta \VUgras{apelsine}. Me lähenesime \VUgras{müüjannale}
\textsuperscript{(94)} ja kui \VUgras{kasvataja} \textsuperscript{(93)}, kes valis oma
kasvandiku jaoks, oli oma ostu lõpetanud, lasksime end teenindada.

%%K t14-22-1 p
22. --- Koju tagasi tulles kohtasime me mitut tuttavat: mu perekonna vana sõbrannat,
kes toetas kätt oma \VUgras{seltsidaamile} \textsuperscript{(95)}, sildade ja teede
\VUgras{inseneri} \textsuperscript{(96)} ja üht oma nõbudest koos oma
\VUgras{lapsehoidjaga} \textsuperscript{(98)}.

%%K t14-22-2 p
Siis kohtasime me oma ema \VUgras{kübarameistrit} \textsuperscript{(97)} ja kaht linnale
võõrast \VUgras{munka} \textsuperscript{(99)}, kes tulid meie juurde, et küsida mu isalt
teed jaama.

\VUsaut{6.0mm}
\VUfilet{22mm}
